%% GENERATED by tools/gen_latex.py -- edit content/ch08.py instead.
%% Build:  latexmk -pdf ch08.tex     (or: pdflatex ch08.tex, twice)
\documentclass[aspectratio=169,11pt,t]{beamer}
\usepackage{itbpro}

\coursetitle{Image Classification}
\coursesubtitle{ConvNets from first principles, then the three strategies that make small image datasets tractable --- from 80\% to 98.6\% on the same 2,000 pictures.}
\coursekicker{Chapter 8}
\coursesource{Chollet \& Watson, Deep Learning with Python 3e --- chapter 8}
\courseduration{3 hours (2 sessions)}
\coursepresenter{Rahman Indra Kesuma, S.Kom., M.Cs. $\cdot$ Prof. Bambang Riyanto Trilaksono}{Teaching Assistant $\cdot$ Lead Instructor}
\coursebrand{AI for Professional \textperiodcentered\ ITB \textperiodcentered\ Ch 8}

\title{Image Classification}
\author{Rahman Indra Kesuma, S.Kom., M.Cs. $\cdot$ Prof. Bambang Riyanto Trilaksono}
\date{}

\begin{document}
\covertitle

\framekicker{Learning outcomes}
\begin{frame}[shrink=25]{Session Objectives}
\leadpar{By the end of this session, participants can:}
\begin{isteps}
\item Explain why a convolution layer beats a dense layer on images: \textbf{translation invariance} and \textbf{spatial hierarchies}.
\item Read a ConvNet summary and predict how \textbf{feature-map shape and depth} change layer by layer.
\item Say what \textbf{max pooling} is for, and what breaks in a ConvNet without it.
\item Build an image pipeline with \textbf{\inlinecode{image\_dataset\_from\_directory}} and \textbf{\inlinecode{tf.data}}, including prefetching.
\item Apply \textbf{data augmentation} and explain what it can and cannot fix.
\item Use a pretrained model two ways --- \textbf{fast feature extraction} and \textbf{end-to-end with augmentation} --- and know the trade-off.
\item \textbf{Fine-tune} a pretrained base in the right order, at the right learning rate, and know which layers to leave frozen.
\end{isteps}

\vskip 6pt
\vskip 4pt
\reslink{Book}{Chapter 8 --- full text}{https://deeplearningwithpython.io/chapters/chapter08\_image-classification}
\reslink{Notebook}{01\_convnet\_on\_mnist.ipynb}{../../notebooks/ch08/01\_convnet\_on\_mnist.ipynb}
\reslink{Notebook}{02\_small\_dataset\_from\_scratch.ipynb}{../../notebooks/ch08/02\_small\_dataset\_from\_scratch.ipynb}
\reslink{Notebook}{03\_data\_augmentation.ipynb}{../../notebooks/ch08/03\_data\_augmentation.ipynb}
\reslink{Notebook}{04\_feature\_extraction\_and\_finetuning.ipynb}{../../notebooks/ch08/04\_feature\_extraction\_and\_finetuning.ipynb}
\reslink{Repo}{Author's official notebook}{https://github.com/fchollet/deep-learning-with-python-notebooks/blob/master/chapter08\_image-classification.ipynb}
\vskip 2pt
\end{frame}

\framekicker{Roadmap}
\begin{frame}[shrink=25]{One problem, three strategies, four numbers}
\begin{minipage}[t]{0.241\textwidth}
\istat{$\approx$80\%}{small ConvNet from scratch, no regularisation}
\end{minipage}%
\hfill
\begin{minipage}[t]{0.241\textwidth}
\istat{$\approx$84\%}{the same, with data augmentation}
\end{minipage}%
\hfill
\begin{minipage}[t]{0.241\textwidth}
\istat{98.1\%}{feature extraction from a pretrained model}
\end{minipage}%
\hfill
\begin{minipage}[t]{0.241\textwidth}
\istat{98.6\%}{fine-tuning that pretrained model}
\end{minipage}%
{\fontsize{9.4}{12.6}\selectfont\color{fg2}All four use the \textbf{same 2,000 training images}. The lesson of the chapter is in the gap between the second number and the third --- \hilite{reusing learned features is worth more than any amount of tuning your own small model}\par}\vskip 4pt
\end{frame}

\coursesection{01}{Introduction to ConvNets}{Why a convolution beats a dense layer on images.}

\framekicker{Section 8.1 $\cdot$ listing 8.1}
\begin{frame}[shrink=25]{A basic ConvNet is a stack of two layer types}
{\fontsize{9.4}{12.6}\selectfont\color{fg2}Before any theory, here is one that works. It classifies MNIST digits, the same task chapter 2 solved with a dense network.\par}\vskip 4pt
\icode{python}{listing 8.1 --- a small ConvNet}{listings/ch08-01.py}
{\fontsize{9.4}{12.6}\selectfont\color{fg2}A ConvNet takes tensors of shape \inlinecode{(image\_height, image\_width, image\_channels)}, not counting the batch dimension.\par}\vskip 4pt
\end{frame}

\framekicker{Section 8.1 $\cdot$ listing 8.2}
\begin{frame}[shrink=25]{What the summary tells you}
{\fontsize{9.4}{12.6}\selectfont\color{fg2}Two patterns are visible in the shapes, and they hold for almost every ConvNet you will meet.\par}\vskip 4pt
\ioutput{listings/ch08-02.txt}
\begin{ibullets}
\item \textbf{Width and height shrink} as you go deeper: 28 $\rightarrow$ 26 $\rightarrow$ 13 $\rightarrow$ 11 $\rightarrow$ 5 $\rightarrow$ 3.
\item \textbf{Depth grows}: 1 $\rightarrow$ 64 $\rightarrow$ 128 $\rightarrow$ 256, controlled by the first argument to \inlinecode{Conv2D}.
\end{ibullets}
\end{frame}

\framekicker{Section 8.1 $\cdot$ listing 8.3--8.4}
\begin{frame}[shrink=25]{The result, against the dense network from chapter 2}
{\fontsize{9.4}{12.6}\selectfont\color{fg2}Same data, same loss, same number of epochs. Only the architecture changed.\par}\vskip 4pt
\icode{python}{listing 8.3 --- training it}{listings/ch08-03.py}
\ioutput{listings/ch08-04.txt}
\begin{minipage}[t]{0.494\textwidth}
\istat{97.8\% $\rightarrow$ 99.1\%}{dense network vs ConvNet}
\end{minipage}%
\hfill
\begin{minipage}[t]{0.494\textwidth}
\istat{$\approx$ 60\%}{relative reduction in error rate}
\end{minipage}%
\end{frame}

\framekicker{Section 8.1.1}
\begin{frame}[shrink=25]{The fundamental difference: global versus local}
\begin{center}
\includegraphics[width=0.92\linewidth,height=0.60\textheight,keepaspectratio]{figs/ch08-localglobal.tex.pdf}
\end{center}
\vskip 3pt{\color{ink3}\fontsize{7.4}{9.4}\selectfont The one architectural change that produced that 60\% error reduction.\par}
\end{frame}

\framekicker{Section 8.1.1}
\begin{frame}[shrink=25]{Property 1 --- translation invariance}
{\fontsize{9.4}{12.6}\selectfont\color{fg2}Having learned a pattern in the lower-right corner, a ConvNet can recognise it \textbf{anywhere} --- upper-left included. A dense network would have to learn it again for the new position.\par}\vskip 4pt
\begin{iband}[signal]
This makes ConvNets \textbf{data efficient} on images, because the visual world is fundamentally translation invariant. They \hilite{need fewer training samples} to learn representations that generalise.
\end{iband}
\end{frame}

\framekicker{Section 8.1.1}
\begin{frame}[shrink=25]{Property 2 --- spatial hierarchies}
\begin{center}
\includegraphics[width=0.92\linewidth,height=0.60\textheight,keepaspectratio]{figs/ch08-hierarchy.tex.pdf}
\end{center}
\vskip 3pt{\color{ink3}\fontsize{7.4}{9.4}\selectfont Elementary lines and textures combine into eyes and ears, which combine into high-level concepts.\par}
{\fontsize{9.4}{12.6}\selectfont\color{fg2}This works because \textbf{the visual world is itself spatially hierarchical}. The architecture matches the structure of the data --- which is exactly what chapter 5 called a good \hilite{architecture prior}\par}\vskip 4pt
\end{frame}

\framekicker{Section 8.1.1}
\begin{frame}[shrink=25]{Feature maps, filters, and response maps}
{\fontsize{9.4}{12.6}\selectfont\color{fg2}Convolutions operate over rank-3 tensors called \textbf{feature maps}: two spatial axes plus a depth axis. For an RGB image the depth is 3; for grayscale it is 1.\par}\vskip 4pt
\begin{ibullets}
\item The output's depth is \textbf{a parameter of the layer}, and its channels no longer stand for colours --- they stand for \textbf{filters}.
\item A filter encodes a specific aspect of the input; at a high level, one filter might encode \textit{"presence of a face"}.
\item In the MNIST example the first layer turns \inlinecode{(28, 28, 1)} into \inlinecode{(26, 26, 64)} --- it computes \textbf{64 filters} over its input.
\item Each of those 64 channels is a 26$\times$26 grid: a \textbf{response map} saying \hilite{how strongly that filter's pattern appears at each location}.
\end{ibullets}
\end{frame}

\framekicker{Section 8.1.1}
\begin{frame}[shrink=25]{How the convolution actually runs}
\begin{center}
\includegraphics[width=0.92\linewidth,height=0.60\textheight,keepaspectratio]{figs/book/figure-8-4.png}
\end{center}
\vskip 3pt{\color{ink3}\fontsize{7.4}{9.4}\selectfont Figure 8.4 --- a window slides over the input, each 3D patch is turned into a vector by the same learned kernel, and the vectors are reassembled into an output map.~\textcolor{ink3}{\footnotesize---~Chollet \& Watson, \emph{Deep Learning with Python}, 3rd ed. (Manning)}\par}
\end{frame}

\framekicker{Section 8.1.1}
\begin{frame}[shrink=25]{One position of that window, computed}
\begin{center}
\includegraphics[width=0.97\linewidth,height=0.56\textheight,keepaspectratio]{figs/ch08-conv-compute.tex.pdf}
\end{center}
\vskip 3pt{\color{ink3}\fontsize{7.4}{9.4}\selectfont A 3$\times$3 kernel over one 3$\times$3 patch. This kernel is a Sobel filter, which responds to vertical edges.\par}
\end{frame}

\framekicker{Section 8.1.1}
\begin{frame}[shrink=25]{Reading the arithmetic: three things it settles}
\begin{isteps}
\item \textbf{A convolution is multiply-and-add, and nothing else.} Nine products and a sum give one output cell. Slide the window to the next position and do it again --- that is the whole operation.
\item \textbf{The kernel is what gets learned.} Its nine numbers are \inlinecode{1, 0, -1 / 2, 0, -2 / 1, 0, -1} here because a person chose a Sobel filter. In a ConvNet \hilite{those nine numbers are parameters}, found by gradient descent, and nobody chooses what they end up detecting.
\item \textbf{The same nine numbers are reused at every position.} That is weight sharing: it is why a ConvNet has so few parameters compared with a \inlinecode{Dense} layer over the same image, and why a pattern learned in one corner is recognised in the other.
\end{isteps}
\begin{iband}[signal]
Why this kernel finds vertical edges: the left column is positive and the right column is negative, so the sum is large wherever \textbf{the left side is brighter than the right}, and near zero on flat ground. Run it on a patch where every cell is equal and the answer is exactly 0.
\end{iband}
\end{frame}

\framekicker{Section 8.1.1}
\begin{frame}[shrink=25]{The two parameters that define a convolution}
\begin{itable}{>{\raggedright\arraybackslash}p{0.282\textwidth}>{\raggedright\arraybackslash}p{0.2068\textwidth}>{\raggedright\arraybackslash}p{0.4512\textwidth}}
\thead{Parameter} & \thead{Typical value} & \thead{What it controls} \\ \midrule
\textbf{Patch size} & 3 $\times$ 3 or 5 $\times$ 5 & How large a neighbourhood each output value can see. \\
\textbf{Output depth} & 32 to 512 & How many filters are computed --- how many different patterns the layer can look for. \\
\end{itable}
\icode{python}{how Keras spells it}{listings/ch08-05.py}
{\fontsize{9.4}{12.6}\selectfont\color{fg2}Because the kernel is shared across all positions, \hilite{parameter count depends on the window and depth, not on the image size}\par}\vskip 4pt
\end{frame}

\framekicker{Section 8.1.2}
\begin{frame}[shrink=25]{Max pooling: downsample, aggressively}
{\fontsize{9.4}{12.6}\selectfont\color{fg2}After each \inlinecode{MaxPooling2D} the feature map halves: 26$\times$26 becomes 13$\times$13. It extracts windows and outputs the maximum value of each channel.\par}\vskip 4pt
\begin{itable}{>{\raggedright\arraybackslash}p{0.2444\textwidth}>{\raggedright\arraybackslash}p{0.3478\textwidth}>{\raggedright\arraybackslash}p{0.3478\textwidth}}
\thead{} & \thead{Convolution} & \thead{Max pooling} \\ \midrule
Transformation & A \textbf{learned} linear transform (the kernel). & A \textbf{hardcoded} max operation. \\
Typical window & 3 $\times$ 3, stride 1 & 2 $\times$ 2, stride 2 \\
Effect on size & Shrinks slightly (no padding). & \textbf{Halves} each spatial dimension. \\
\end{itable}
\end{frame}

\framekicker{Section 8.1.2 $\cdot$ listing 8.5}
\begin{frame}[shrink=25]{What happens without it}
{\fontsize{9.4}{12.6}\selectfont\color{fg2}The book builds the broken version deliberately: the same network with every pooling layer removed.\par}\vskip 4pt
\icode{python}{listing 8.5 --- an incorrectly structured ConvNet}{listings/ch08-06.py}
\begin{iband}[rose]
\textbf{It is not conducive to learning a spatial hierarchy.} The 3$\times$3 windows in the third layer still only see information from \hilite{7$\times$7 windows of the original image} --- far too small a view to assemble a whole digit.
\end{iband}
\end{frame}

\coursesection{02}{Training from scratch on a small dataset}{2,000 pictures of cats and dogs. A situation you will actually meet.}

\framekicker{Section 8.2}
\begin{frame}[shrink=25]{The setup, and why it is realistic}
{\fontsize{9.4}{12.6}\selectfont\color{fg2}Having to train an image classifier on very little data is \textbf{a common situation} in professional computer vision. A "few" samples can mean a few hundred to a few tens of thousands.\par}\vskip 4pt
\begin{minipage}[t]{0.3253\textwidth}
\istat{2,000}{training images}
\end{minipage}%
\hfill
\begin{minipage}[t]{0.3253\textwidth}
\istat{1,000}{validation images}
\end{minipage}%
\hfill
\begin{minipage}[t]{0.3253\textwidth}
\istat{2,000}{test images}
\end{minipage}%
{\fontsize{9.4}{12.6}\selectfont\color{fg2}Drawn from the original Kaggle Dogs vs. Cats dataset, which was released as a competition in late 2013 --- \hilite{before ConvNets were mainstream}\par}\vskip 4pt
\end{frame}

\framekicker{Section 8.2.1}
\begin{frame}[shrink=25]{Why deep learning still applies with little data}
\begin{minipage}[t]{0.494\textwidth}
\begin{icardgood}
\cardh{ConvNets are data efficient}
Because they learn \textbf{local, translation-invariant} features, they give reasonable results on small image datasets without any custom feature engineering.
\end{icardgood}
\end{minipage}%
\hfill
\begin{minipage}[t]{0.494\textwidth}
\begin{icardgood}
\cardh{Models are repurposable}
A classifier trained on a large dataset can be reused on a significantly different problem with minor changes. \textbf{Feature reuse is one of deep learning's greatest strengths.}
\end{icardgood}
\end{minipage}%
{\fontsize{9.4}{12.6}\selectfont\color{fg2}What counts as \textit{enough} is relative --- to the size and depth of the model. A few tens of samples will not do; a few hundred can suffice if the model is small, well regularised, and the task is simple.\par}\vskip 4pt
\end{frame}

\framekicker{Section 8.2.3 $\cdot$ listing 8.7}
\begin{frame}[shrink=25]{A bigger ConvNet for bigger images}
{\fontsize{9.4}{12.6}\selectfont\color{fg2}Same alternating pattern as before, with two extra stages --- the images are larger and the problem harder. Note the \inlinecode{Rescaling} layer, which moves the preprocessing \textbf{inside the model}.\par}\vskip 4pt
\icode{python}{listing 8.7 --- dogs vs cats, from scratch}{listings/ch08-07.py}
{\fontsize{9.4}{12.6}\selectfont\color{fg2}Putting \inlinecode{Rescaling} in the model means the deployed artefact carries its own preprocessing --- one fewer thing to get wrong in production.\par}\vskip 4pt
\end{frame}

\framekicker{Section 8.2.3}
\begin{frame}[shrink=25]{The pattern you will see in almost every ConvNet}
\begin{center}
\includegraphics[width=0.92\linewidth,height=0.60\textheight,keepaspectratio]{figs/ch08-shapes.tex.pdf}
\end{center}
\vskip 3pt{\color{ink3}\fontsize{7.4}{9.4}\selectfont Depth rises from 32 to 512 while the spatial size falls from 180$\times$180 to 7$\times$7.\par}
\begin{iband}[signal]
\textbf{Feature-map depth increases while feature-map size decreases.} Information moves out of \textit{where} and into \textit{what} --- \hilite{and that is the whole shape of a classification ConvNet}
\end{iband}
\end{frame}

\framekicker{Section 8.2.4}
\begin{frame}[shrink=25]{From JPEG files to batched tensors}
{\fontsize{9.4}{12.6}\selectfont\color{fg2}The data sits on disk as JPEGs. Five steps are needed to make it model-ready, and Keras does all of them in one call.\par}\vskip 4pt
\begin{isteps}
\item Read the picture files.
\item Decode the JPEG content into RGB grids of pixels.
\item Convert those into floating-point tensors.
\item Resize them to a shared size --- here 180 $\times$ 180.
\item Pack them into batches.
\end{isteps}
\end{frame}

\framekicker{Section 8.2.4 $\cdot$ listing 8.9}
\begin{frame}[shrink=25]{\inlinecode{image\_dataset\_from\_directory} does all five}
{\fontsize{9.4}{12.6}\selectfont\color{fg2}It lists the subdirectories, treats each as a class, indexes the files, and returns a \inlinecode{tf.data.Dataset} that shuffles, decodes, resizes, and batches.\par}\vskip 4pt
\icode{python}{listing 8.9 --- building the pipeline}{listings/ch08-08.py}
{\fontsize{9.4}{12.6}\selectfont\color{fg2}The directory layout \textbf{is} the label set: one subfolder per class, and the folder names become the classes.\par}\vskip 4pt
\end{frame}

\framekicker{Section 8.2.4}
\begin{frame}[shrink=25]{What a \inlinecode{tf.data.Dataset} buys you}
\begin{ibullets}
\item It works with \textbf{any backend} --- JAX and PyTorch included --- not only TensorFlow.
\item It is an \textbf{iterator}: usable in a \inlinecode{for} loop, and passable straight to \inlinecode{fit()}.
\item It \textbf{parallelises preprocessing across CPU cores}.
\item It does \textbf{asynchronous prefetching}: preparing the next batch while the model is still working on the previous one, so \hilite{execution never stalls}.
\end{ibullets}
\begin{iband}[signal]
Those last two are the difference between a GPU that is busy and a GPU that spends most of its time waiting for JPEGs to decode.
\end{iband}
\end{frame}

\coursesection{03}{Data augmentation}{Making more training data out of the training data you have.}

\framekicker{Section 8.2.5}
\begin{frame}[shrink=25]{The idea, and where the layers go}
{\fontsize{9.4}{12.6}\selectfont\color{fg2}Overfitting comes from having too few samples. Augmentation generates more of them by applying \textbf{random but believable} transformations, so the model \hilite{never sees exactly the same picture twice}\par}\vskip 4pt
\begin{itable}{>{\raggedright\arraybackslash}p{0.282\textwidth}>{\raggedright\arraybackslash}p{0.188\textwidth}>{\raggedright\arraybackslash}p{0.47\textwidth}}
\thead{Where you put the layers} & \thead{Runs on} & \thead{Trade-off} \\ \midrule
\textbf{Inside the model}, before \inlinecode{Rescaling} & GPU & Sometimes faster, but the augmentation ships with the model. \\
\textbf{In the data pipeline}, via \inlinecode{map()} & CPU, in parallel & \textbf{Usually the better option} --- and the one the book takes. \\
\end{itable}
\end{frame}

\framekicker{Section 8.2.5 $\cdot$ listing 8.18}
\begin{frame}[shrink=25]{Three transformations, and prefetching}
{\fontsize{9.4}{12.6}\selectfont\color{fg2}The layers are defined as a list, wrapped in a function, and mapped over the dataset.\par}\vskip 4pt
\icode{python}{listing 8.18 --- the augmentation stage}{listings/ch08-09.py}
\begin{ibullets}
\item \inlinecode{RandomFlip("horizontal")} --- flips a random \textbf{50\%} of images.
\item \inlinecode{RandomRotation(0.1)} --- rotates by a random value in \textbf{[$-$10\%, +10\%]} of a full circle, i.e. \textbf{$\pm$36 degrees}.
\item \inlinecode{RandomZoom(0.2)} --- zooms in or out by a random factor in \textbf{[$-$20\%, +20\%]}.
\end{ibullets}
\end{frame}

\framekicker{Section 8.2.5}
\begin{frame}[shrink=25]{What augmentation cannot do}
\begin{iband}[amber]
The model will never see the same input twice --- but the inputs it does see are \textbf{still heavily intercorrelated}, because they come from a small number of originals. \hilite{You cannot produce new information; you can only remix existing information.}
\end{iband}
\begin{minipage}[t]{0.494\textwidth}
\istat{$\approx$ 80\%}{from scratch, no augmentation}
\end{minipage}%
\hfill
\begin{minipage}[t]{0.494\textwidth}
\istat{$\approx$ 84\%}{from scratch, with augmentation}
\end{minipage}%
{\fontsize{9.4}{12.6}\selectfont\color{fg2}Four points for a change that costs nothing but CPU time --- worth having, and clearly not enough on its own.\par}\vskip 4pt
\end{frame}

\coursesection{04}{Using a pretrained model}{Where the remaining fourteen points come from.}

\framekicker{Section 8.3}
\begin{frame}[shrink=25]{Why features transfer between problems}
{\fontsize{9.4}{12.6}\selectfont\color{fg2}A model trained on a large, general dataset learns a spatial hierarchy of features that acts as \textbf{a generic model of the visual world} --- useful even for classes it never saw.\par}\vskip 4pt
{\fontsize{9.4}{12.6}\selectfont\color{fg2}The book uses \textbf{Xception} trained on ImageNet: 1.4 million labelled images across 1,000 classes. ImageNet contains many cat and dog breeds, so it should transfer well here.\par}\vskip 4pt
\begin{iband}[signal]
Such portability across problems is \textbf{a key advantage of deep learning over older shallow methods}, and it is what makes it effective on small-data problems at all.
\end{iband}
\end{frame}

\framekicker{Section 8.3}
\begin{frame}[shrink=25]{Two ways to reuse it}
\begin{center}
\includegraphics[width=0.92\linewidth,height=0.60\textheight,keepaspectratio]{figs/ch08-transfer.tex.pdf}
\end{center}
\vskip 3pt{\color{ink3}\fontsize{7.4}{9.4}\selectfont Feature extraction swaps the classifier; fine-tuning then also adjusts the top of the base.\par}
\end{frame}

\framekicker{Section 8.3.1}
\begin{frame}[shrink=25]{Why only the convolutional base is reused}
{\fontsize{9.4}{12.6}\selectfont\color{fg2}A ConvNet has two parts: a \textbf{convolutional base} (the pooling and convolution layers, also called the \textit{backbone}) and a densely connected classifier. Only the first is worth keeping.\par}\vskip 4pt
\begin{ibullets}
\item The base produces \textbf{presence maps of generic concepts} --- useful regardless of the vision problem.
\item The classifier's representations are \textbf{specific to the classes it was trained on}.
\item Worse, dense layers \textbf{discard spatial information}: they no longer know \textit{where} anything is. For problems where location matters they are \hilite{largely useless}.
\end{ibullets}
\end{frame}

\framekicker{Section 8.3.1}
\begin{frame}[shrink=25]{How much of the base to reuse}
{\fontsize{9.4}{12.6}\selectfont\color{fg2}Generality depends on depth. \textbf{Earlier layers extract local, highly generic feature maps} --- edges, colours, textures. \textbf{Higher layers extract abstract concepts} such as \textit{cat ear} or \textit{dog eye}.\par}\vskip 4pt
\begin{iband}[amber]
So if your new dataset differs a lot from the original, you may be better off using \hilite{only the first few layers} for feature extraction rather than the entire base.
\end{iband}
\end{frame}

\framekicker{Section 8.3.1 $\cdot$ listing 8.23--8.24}
\begin{frame}[shrink=25]{Loading a pretrained backbone}
{\fontsize{9.4}{12.6}\selectfont\color{fg2}KerasHub ships Keras implementations of popular architectures with downloadable weights. Note that the preprocessing comes \textbf{paired with the checkpoint}.\par}\vskip 4pt
\icode{python}{listing 8.23--8.24 --- backbone and preprocessing}{listings/ch08-10.py}
\begin{ibullets}
\item \textbf{backbone} is KerasHub's word for the feature extractor without its classification head.
\item The \textbf{41} is a naming convention: 41 trainable layers. Pretrained ConvNets are conventionally named by depth.
\item \inlinecode{ImageConverter} matters more than it looks: every pretrained ConvNet expects a particular input range, and feeding the wrong one forces the model to \hilite{relearn how to see}.
\end{ibullets}
\end{frame}

\framekicker{Section 8.3.1}
\begin{frame}[shrink=25]{Fast, or augmentable --- pick one}
\begin{center}
\includegraphics[width=0.92\linewidth,height=0.60\textheight,keepaspectratio]{figs/ch08-twoways.tex.pdf}
\end{center}
\vskip 3pt{\color{ink3}\fontsize{7.4}{9.4}\selectfont The base is by far the most expensive part of the pipeline, so running it once versus every epoch is the whole trade-off.\par}
\end{frame}

\framekicker{Section 8.3.1 $\cdot$ listing 8.25}
\begin{frame}[shrink=25]{Fast route --- cache the features once}
{\fontsize{9.4}{12.6}\selectfont\color{fg2}Run the frozen base over each dataset a single time and keep the outputs as NumPy arrays.\par}\vskip 4pt
\icode{python}{listing 8.25 --- extracting features}{listings/ch08-11.py}
\ioutput{listings/ch08-12.txt}
\end{frame}

\framekicker{Section 8.3.1 $\cdot$ listing 8.26}
\begin{frame}[shrink=25]{\ldots{}then train a tiny classifier on the cache}
{\fontsize{9.4}{12.6}\selectfont\color{fg2}The classifier now sees 6$\times$6$\times$2048 feature maps rather than images, so it is small and trains in seconds.\par}\vskip 4pt
\icode{python}{listing 8.26 --- the dense classifier}{listings/ch08-13.py}
\ioutput{listings/ch08-14.txt}
\end{frame}

\framekicker{Section 8.3.1}
\begin{frame}[shrink=25]{Read that result carefully}
\begin{iband}[amber]
It is \textbf{a slightly unfair comparison}: ImageNet contains many dog and cat instances, so the pretrained model already holds \hilite{the exact knowledge this task needs}. That will not always be true.
\end{iband}
{\fontsize{9.4}{12.6}\selectfont\color{fg2}The curves also show \textbf{overfitting almost from the start}, despite a fairly large dropout rate --- because this route cannot use augmentation, which is essential on small image datasets.\par}\vskip 4pt
\end{frame}

\framekicker{Section 8.3.1 $\cdot$ listing 8.28--8.29}
\begin{frame}[shrink=25]{Slow route --- freeze the base and chain it}
{\fontsize{9.4}{12.6}\selectfont\color{fg2}Freezing means preventing weights from updating. Without it, the randomly initialised classifier would send huge updates backwards and \hilite{destroy the representations the base already learned}\par}\vskip 4pt
\icode{python}{listing 8.28--8.29 --- freezing}{listings/ch08-15.py}
\ioutput{listings/ch08-16.txt}
{\fontsize{9.4}{12.6}\selectfont\color{fg2}Setting \inlinecode{trainable = False} \textbf{empties the list of trainable weights} --- which is exactly how you verify that a freeze took effect.\par}\vskip 4pt
\end{frame}

\framekicker{Section 8.3.1}
\begin{frame}[shrink=25]{The chained model, and a recompile trap}
{\fontsize{9.4}{12.6}\selectfont\color{fg2}Now the preprocessing, the frozen base, and a new classifier are one model --- so augmentation applies to every image on every epoch.\par}\vskip 4pt
\icode{python}{the end-to-end model}{listings/ch08-17.py}
\begin{iband}[rose]
Only the two \inlinecode{Dense} layers train --- four weight tensors in total. And a trap worth memorising: \textbf{if you change trainability after compiling, you must recompile}, or \hilite{the change is silently ignored}
\end{iband}
\end{frame}

\coursesection{05}{Fine-tuning}{Unfreezing the top of the base --- carefully, and last.}

\framekicker{Section 8.3.2}
\begin{frame}[shrink=25]{The five steps, and why the order is fixed}
\begin{center}
\includegraphics[width=0.92\linewidth,height=0.60\textheight,keepaspectratio]{figs/ch08-finetune.tex.pdf}
\end{center}
\vskip 3pt{\color{ink3}\fontsize{7.4}{9.4}\selectfont Steps 1--3 are exactly the feature-extraction workflow. Fine-tuning is what you do afterwards.\par}
\begin{iband}[rose]
The base can only be fine-tuned \textbf{once the classifier on top is already trained}. Otherwise the error signal propagating back is too large, and \hilite{the representations being fine-tuned are destroyed}
\end{iband}
\end{frame}

\framekicker{Section 8.3.2}
\begin{frame}[shrink=25]{Partial fine-tuning, and why the top layers}
{\fontsize{9.4}{12.6}\selectfont\color{fg2}With large pretrained models it is common to unfreeze only the top few layers of the base.\par}\vskip 4pt
\begin{minipage}[t]{0.485\textwidth}
\begin{ibullets}
\item Earlier layers encode \textbf{generic, reusable} features; higher ones encode \textbf{specialised} features.
\item It is the specialised ones that need repurposing, so there are \hilite{fast-decreasing returns} to fine-tuning lower layers.
\item The base has \textbf{15 million parameters} --- training all of them on 2,000 images risks overfitting.
\end{ibullets}

\end{minipage}%
\hfill
\begin{minipage}[t]{0.485\textwidth}
\icode{python}{unfreezing the top four}{listings/ch08-18.py}

\end{minipage}%
\begin{iband}[amber]
One layer type to leave alone: \textbf{do not unfreeze \inlinecode{BatchNormalization} layers}. Chapter 9 explains why.
\end{iband}
\end{frame}

\framekicker{Section 8.3.2 $\cdot$ listing 8.31}
\begin{frame}[shrink=25]{Fine-tune at a very low learning rate}
{\fontsize{9.4}{12.6}\selectfont\color{fg2}The learning rate is deliberately tiny --- the goal is to \textbf{limit the magnitude of the changes} made to representations that are already good.\par}\vskip 4pt
\icode{python}{listing 8.31 --- fine-tuning}{listings/ch08-19.py}
\begin{iband}[signal]
Updates that are too large \hilite{harm the representations you are trying to keep}. This is one of the few places in the book where the learning rate is prescribed rather than searched for.
\end{iband}
\end{frame}

\framekicker{Section 8.3.2}
\begin{frame}[shrink=25]{The final number, in context}
{\fontsize{9.4}{12.6}\selectfont\color{fg2}The checkpoint callback kept the best version, so the model that gets evaluated is the best one seen during the run rather than the last one.\par}\vskip 4pt
\icode{python}{evaluating the fine-tuned model}{listings/ch08-20.py}
\ioutput{listings/ch08-21.txt}
{\fontsize{9.4}{12.6}\selectfont\color{fg2}In the original Kaggle competition this \textbf{would have been among the top results} --- though not a fair comparison, since the pretrained features already carried prior knowledge about cats and dogs that competitors could not use.\par}\vskip 4pt
\begin{iband}[signal]
The fairer point: this was reached using \textbf{about 10\% of the training data} available in that competition. \hilite{There is a huge difference between training on 20,000 samples and on 2,000.}
\end{iband}
\end{frame}

\framekicker{Summary}
\begin{frame}[shrink=25]{What has to survive this chapter}
\begin{isteps}
\item \textbf{Convolutions learn local patterns}, which buys translation invariance and spatial hierarchies --- and makes ConvNets data efficient on images.
\item \textbf{Depth rises as spatial size falls.} That shape is nearly universal in classification ConvNets.
\item \textbf{Max pooling is not optional}: without it, deep layers still only see a tiny window of the original image.
\item \textbf{\inlinecode{image\_dataset\_from\_directory} + \inlinecode{tf.data}} turns a folder of JPEGs into a parallel, prefetching pipeline.
\item \textbf{Augmentation remixes information; it cannot create it.} Worth about four points here.
\end{isteps}
\end{frame}

\framekicker{Summary}
\begin{frame}[shrink=25]{\ldots{}and the part that changes how you plan a project}
\begin{minipage}[t]{0.494\textwidth}
\begin{icardgood}
\cardh{Reuse beats tuning}
Feature extraction from a pretrained base took the same 2,000 images from \textbf{84\% to 98.1\%} --- more than anything you could do to a small model of your own.
\end{icardgood}
\end{minipage}%
\hfill
\begin{minipage}[t]{0.494\textwidth}
\begin{icardaccent}
\cardh{Fine-tune last, and gently}
Only after the new classifier is trained, at a \textbf{very low learning rate}, and with \inlinecode{BatchNormalization} left frozen.
\end{icardaccent}
\end{minipage}%
\begin{iband}[signal]
Which is why the first question on a new vision problem is not \textit{what architecture?} but \hilite{\textit{whose features can I start from?}}
\end{iband}
\vskip 4pt
\reslink{NOTEBOOK}{04\_feature\_extraction\_and\_finetuning.ipynb}{../../course-slides/notebooks/ch08/04\_feature\_extraction\_and\_finetuning.ipynb}
\reslink{NEXT}{Chapter 9 --- ConvNet architecture patterns}{../ch09/index.html}
\vskip 2pt
\end{frame}

\end{document}
